% GENERATED FILE. Edit canonical lessons, then run npm run generate:books.
% canonical-source-hash: fnv1a64:b46b1fce38f17472
% canonical-lessons: PT-C339-acordo, PT-C339-concordar, PT-C339-admitir, PT-C339-consenso, PT-C339-ceder

\chapter{Chegar a acordo}
\label{ch:pt-chegar-a-acordo}
\hlchaptermodality{339}

\begin{chapteropening}
\textbf{By the end of this chapter:} \emph{I can reach an agreement: agree with someone, admit they are right, find a consensus and give way.}

Complete the last lesson of chapter 339: I can reach an agreement: agree with someone, admit they are right, find a consensus and give way.
\end{chapteropening}

\section[o acordo]{o acordo — the agreement, the deal}

\label{lesson:PT-C339-acordo}



\begin{quote}
Before the new one: say the Portuguese for the objection, then the Portuguese for to contradict.
\end{quote}

\subsection*{You'll want to know: o acordo}
\textbf{o acordo} — \textquotedblleft{}the agreement, the deal\textquotedblright{}.

\begin{quote}
Depois de duas horas, chegámos a um acordo.
\end{quote}

\begin{quote}
\emph{After two hours, we reached an agreement.}
\end{quote}

\emph{Estar de acordo} — to agree; \emph{de acordo!} — agreed, OK.

\begin{cousinweb}
\emph{acordo} is built on the verb \emph{acordar} in its old sense \textquotedblleft{}to agree\textquotedblright{}, usually traced to Vulgar Latin \emph{accordāre}, a formation on Latin \emph{cor, cordis} \textquotedblleft{}heart\textquotedblright{} (Lewis \& Short): hearts that beat together.
\end{cousinweb}

\begin{grammarlens}[title={What you've built}]
One more word for telling, explaining and describing at B1.
\end{grammarlens}

\subsection*{Your turn}
\begin{itemize}
\raggedright
  \item \emph{Say it:} \emph{o acordo}
  \item \emph{Say it:} \emph{o acordo}, once more
  \item \emph{Say it:} a sentence of your own with \emph{o acordo}
  \item \emph{Recall:} say the Portuguese for the objection, then the Portuguese for to contradict, then say \emph{o acordo} again
\end{itemize}

\begin{tcolorbox}[breakable,colback=teal!4,colframe=teal!35!black,title={Before you move on}]
What does o acordo mean? (\textquotedblleft{}The agreement, the deal\textquotedblright{}.) Say it once more.
\end{tcolorbox}
\section[concordar]{concordar — to agree}

\label{lesson:PT-C339-concordar}



\begin{quote}
Before the new one: say the Portuguese for to contradict, then the Portuguese for the agreement.
\end{quote}

\subsection*{You'll want to know: concordar}
\textbf{concordar} — \textquotedblleft{}to agree\textquotedblright{}.

\begin{quote}
Concordo contigo: é melhor sair cedo.
\end{quote}

\begin{quote}
\emph{I agree with you: it is better to leave early.}
\end{quote}

\emph{Concordar com alguém} — to agree with someone; \emph{concordar em} + infinitive — to agree to do something.

\begin{cousinweb}
\emph{concordar} is a learned borrowing from Latin \emph{concordāre} \textquotedblleft{}to agree together, to be of one mind\textquotedblright{} (EtymDB, Lewis \& Short), from \emph{concors}, built from \emph{con-} \textquotedblleft{}together\textquotedblright{} and \emph{cor, cordis} \textquotedblleft{}heart\textquotedblright{}.
\end{cousinweb}

\begin{grammarlens}[title={What you've built}]
One more word for telling, explaining and describing at B1.
\end{grammarlens}

\subsection*{Your turn}
\begin{itemize}
\raggedright
  \item \emph{Say it:} \emph{concordar}
  \item \emph{Say it:} \emph{concordar}, once more
  \item \emph{Say it:} a sentence of your own with \emph{concordar}
  \item \emph{Recall:} say the Portuguese for to contradict, then the Portuguese for the agreement, then say \emph{concordar} again
\end{itemize}

\begin{tcolorbox}[breakable,colback=teal!4,colframe=teal!35!black,title={Before you move on}]
What does concordar mean? (\textquotedblleft{}To agree\textquotedblright{}.) Say it once more.
\end{tcolorbox}
\section[admitir]{admitir — to admit; to let in}

\label{lesson:PT-C339-admitir}



\begin{quote}
Before the new one: say the Portuguese for the agreement, then the Portuguese for to agree.
\end{quote}

\subsection*{You'll want to know: admitir}
\textbf{admitir} — \textquotedblleft{}to admit; to let in\textquotedblright{}.

\begin{quote}
Tenho de admitir que o teu plano era melhor.
\end{quote}

\begin{quote}
\emph{I have to admit that your plan was better.}
\end{quote}

It also means to let in or accept: \emph{admitir novos sócios} — to take in new members.

\begin{cousinweb}
\emph{admitir} is a learned borrowing from Latin \emph{admittere} \textquotedblleft{}to let in, to allow\textquotedblright{} (EtymDB), built from \emph{ad} \textquotedblleft{}to\textquotedblright{} and \emph{mittere} \textquotedblleft{}to send\textquotedblright{}.
\end{cousinweb}

\begin{grammarlens}[title={What you've built}]
One more word for telling, explaining and describing at B1.
\end{grammarlens}

\subsection*{Your turn}
\begin{itemize}
\raggedright
  \item \emph{Say it:} \emph{admitir}
  \item \emph{Say it:} \emph{admitir}, once more
  \item \emph{Say it:} a sentence of your own with \emph{admitir}
  \item \emph{Recall:} say the Portuguese for the agreement, then the Portuguese for to agree, then say \emph{admitir} again
\end{itemize}

\begin{tcolorbox}[breakable,colback=teal!4,colframe=teal!35!black,title={Before you move on}]
What does admitir mean? (\textquotedblleft{}To admit; to let in\textquotedblright{}.) Say it once more.
\end{tcolorbox}
\section[o consenso]{o consenso — the consensus, the general agreement}

\label{lesson:PT-C339-consenso}



\begin{quote}
Before the new one: say the Portuguese for to agree, then the Portuguese for to admit.
\end{quote}

\subsection*{You'll want to know: o consenso}
\textbf{o consenso} — \textquotedblleft{}the consensus, the general agreement\textquotedblright{}.

\begin{quote}
Não houve consenso na reunião.
\end{quote}

\begin{quote}
\emph{There was no consensus at the meeting.}
\end{quote}

\emph{Chegar a um consenso} — to reach a consensus.

\begin{cousinweb}
\emph{consenso} is a learned borrowing from Latin \emph{cōnsēnsus} \textquotedblleft{}agreement, accordance, unanimity\textquotedblright{} (Lewis \& Short), from \emph{cōnsentīre}, built from \emph{con-} \textquotedblleft{}together\textquotedblright{} and \emph{sentīre} \textquotedblleft{}to feel\textquotedblright{}.
\end{cousinweb}

\begin{grammarlens}[title={What you've built}]
One more word for telling, explaining and describing at B1.
\end{grammarlens}

\subsection*{Your turn}
\begin{itemize}
\raggedright
  \item \emph{Say it:} \emph{o consenso}
  \item \emph{Say it:} \emph{o consenso}, once more
  \item \emph{Say it:} a sentence of your own with \emph{o consenso}
  \item \emph{Recall:} say the Portuguese for to agree, then the Portuguese for to admit, then say \emph{o consenso} again
\end{itemize}

\begin{tcolorbox}[breakable,colback=teal!4,colframe=teal!35!black,title={Before you move on}]
What does o consenso mean? (\textquotedblleft{}The consensus, the general agreement\textquotedblright{}.) Say it once more.
\end{tcolorbox}
\section[ceder]{ceder — to give in, to give way; to hand over}

\label{lesson:PT-C339-ceder}



\begin{quote}
Before the new one: say the Portuguese for to admit, then the Portuguese for the consensus.
\end{quote}

\subsection*{You'll want to know: ceder}
\textbf{ceder} — \textquotedblleft{}to give in, to give way; to hand over\textquotedblright{}.

\begin{quote}
No fim, o meu pai cedeu e deixou-me ir à festa.
\end{quote}

\begin{quote}
\emph{In the end my father gave in and let me go to the party.}
\end{quote}

\emph{Ceder o lugar} — to give up your seat.

\begin{cousinweb}
\emph{ceder} is a learned borrowing from Latin \emph{cēdere} \textquotedblleft{}to go, to withdraw, to yield\textquotedblright{} (EtymDB, Lewis \& Short).
\end{cousinweb}

\begin{grammarlens}[title={What you've built}]
One more word for telling, explaining and describing at B1.
\end{grammarlens}

\subsection*{Your turn}
\begin{itemize}
\raggedright
  \item \emph{Say it:} \emph{ceder}
  \item \emph{Say it:} \emph{ceder}, once more
  \item \emph{Say it:} a sentence of your own with \emph{ceder}
  \item \emph{Recall:} say the Portuguese for to admit, then the Portuguese for the consensus, then say \emph{ceder} again
\end{itemize}

\begin{tcolorbox}[breakable,colback=teal!4,colframe=teal!35!black,title={Before you move on}]
What does ceder mean? (\textquotedblleft{}To give in, to give way; to hand over\textquotedblright{}.) Say it once more.
\end{tcolorbox}
